\chapter{Supplementary Model Comparisons}
\label{app:secondary-results}

\section{Fixed probability combination and training-seed repeat}

For the probability combination in Equation~\ref{eq:mixture}, 144
executions were attempted and 135 were eligible, yielding 1,427
natural rows. The models and mixing rule were frozen before
collection. Table~\ref{tab:mix} compares them on the same
eligible inputs. The five seeds were 53, 67, 79, 97, and 113.
All rows in the two A+C layouts had observed targets, so
conditional mixing and an unconditional half-average produced the
same predictions in those layouts, and the improvement there is that
of the combined system rather than of the conditioning rule.

\begin{table}[tb]
\centering
\caption{Fixed probability combination on newly collected runs from the
same laboratory. All values are execution-weighted NLL. Rows with no
usable target receive the baseline distribution.}
\label{tab:mix}
\begin{tabularx}{\textwidth}{@{}Xrrrr@{}}
\toprule
Condition & Runs / rows & Baseline & Full Step (dual history) & Conditional average \\
\midrule
A+C, flat & 31 / 341 & 0.135230 & 0.107396 & \textbf{0.080156} \\
A+C, nested & 29 / 319 & 0.160900 & \textbf{0.075783} & 0.077815 \\
C-only & 27 / 297 & \textbf{0.060420} & 0.091372 & 0.062979 \\
B+C & 32 / 352 & 0.052914 & 0.059751 & \textbf{0.039371} \\
Original 18-action task & 16 / 118 & 0.645963 & 0.668895 & \textbf{0.640501} \\
\bottomrule
\end{tabularx}
\end{table}

For the full eligible set, the average's loss differences versus the
baseline were $-0.055074$ in the flat A+C layout and $-0.083086$ in the
nested layout. In paired bootstrap estimates using complete
configurations, the Bonferroni-adjusted 97.5\% intervals were
$[-0.092701,-0.027024]$ and $[-0.145118,-0.038204]$,
respectively. The intervals used 30 flat and 26 nested runs from
configurations with both repetitions. The system improved in
four of five seeds on each A+C layout. The conditional rule
retained the baseline distribution on all 96 target-free rows of
the original-task subset. C-only shows the trade-off of the rule:
the conditional average scored 0.062979, compared with 0.060420
for the baseline and 0.054544 for the unconditional average.

A matched retraining comparison fitted all models from scratch with
five new seeds (127, 149, 163, 181, and 199) and reused the same 135
evaluation executions. Table~\ref{tab:matched} shows the resulting
A+C scores. The dual-history average reduced NLL versus its own
baseline in both layouts. Relative to the wider single-history
and BiTA combinations, all four adjusted 98.75\% intervals included
zero. The separate, post-hoc comparison against an average of two
independent baseline checkpoints also yielded intervals that
included zero in both layouts. Because these intervals cross zero,
the gains of the fixed combination cannot be attributed to the
dual-history architecture itself.

\begin{table}[tb]
\centering
\caption{A+C comparison after training with five new seeds. All
combinations apply the same target-conditioned 50:50 average and share
the same baseline within each seed. The evaluation executions are
reused from Table~\ref{tab:mix}.}
\label{tab:matched}
\begin{tabularx}{\textwidth}{@{}Xrr@{}}
\toprule
Predictor & A+C, flat & A+C, nested \\
\midrule
Single baseline & 0.192507 & 0.176440 \\
Baseline + dual history & 0.123421 & \textbf{0.108952} \\
Baseline + wider single history & \textbf{0.120299} & 0.113041 \\
Baseline + BiTA-inspired aggregator & 0.124414 & 0.129035 \\
\bottomrule
\end{tabularx}
\end{table}

\section{Observed-history support on target-poor prefixes}

The frequency table in Equation~\ref{eq:table} was constructed
from 603 training executions and 4,824 natural rows. For the
64-execution, 472-row original-action cohort, the stand-alone
table scored 0.495788, versus 2.843204 for the training prior alone.
For 386 target-free rows in this cohort, mixing the table into
a frozen model yielded 0.313385 overall, compared with 0.317622
for the same mixture using the training prior. The adjusted
98.75\% interval for the difference between the mixtures was
$[-0.026307,0.021147]$. The table thus provides a useful simple
forecast in this cohort, with uncertain additional value once
the learned model and prior mixing are available.

The same key definition was checked on the 660 natural rows across
the two A+C layouts in Table~\ref{tab:mix}. None had an
exact key supported in training, so the table returned the same
training prior on every row. Preserving unclassified START
records in the key yielded the same fallback pattern. A coarser
stage key found support on 161 of 341 flat rows and 148 of 319
nested rows. It scored 1.181161 and 1.268327, respectively. Exact lookup thus
loses its support when benign activity changes the sequence of
received events on which the key depends.
